% 由 work/gen_notes.py 自动生成，请勿手改。
\section{译校批注一览}\label{app:notes}

\noindent 全文共 \textbf{68} 条译校批注。每条的原始位置在对应段落之后的 LaTeX 注释里（\verb|% [Tutor Note: …]|），不进入排版结果；本附录按节次汇总，便于对照原文复查。

\subsection*{第 1 节 引言}
\begin{enumerate}
  \item 本节是全文的路标，一句话概括：把 Fukaya 范畴在余切丛 $T^*X$ 上取半经典极限 $\epsilon\to0$，全纯圆盘的计数就退化为梯度流树的计数；本文补上“对象到自身”的态射，并给出 $A_\infty$ 关系式的两个非全纯证明。\hfill{\footnotesize （源码第 6 行）}
  \item “重标度 $F_a\to\epsilon F_a$”是全篇第一个易错点：$\epsilon\to0$ 时 $L_a=\mathrm{graph}(\epsilon\,\mathrm{d}F_a)$ 退化到零截面，交点 $P_{\alpha,ab}$ 趋于 $\mathrm{d}(F_a-F_b)=0$ 的临界点，全纯圆盘的半径随之趋于零。作者在此省略了“一列全纯圆盘收敛到梯度流树”的紧性论证，标准做法是用 Gromov 紧性与 $T^*X$ 上近全纯曲线的局部描述，见 Fukaya--Oh \cite{fukayaoh} 与 Seidel \cite{seidel} 第 9 章。\hfill{\footnotesize （源码第 26 行）}
  \item “两个证明都不涉及全纯圆盘”是全文纲领，值得停一下。其一只用 BV 形式的纤维积分与树形求和（等价于 Kontsevich--Soibelman 同伦转移定理，见 \cite{kontsevichsoibelman}）；其二只用度量树模空间上一个闭微分形式与 Stokes 定理。换言之，作者把“辛几何”这一层完全剥离，只留下“生成函数之差是 Morse--Smale 函数”与“梯度流”两件事。这也解释了为什么证明能同时给出整系数版本与一般的环系数版本。\hfill{\footnotesize （源码第 30 行）}
  \item (B) 的逻辑方向容易说反。BV 主方程 $\{S_{\mathrm{eff}},S_{\mathrm{eff}}\}=0$ 直接给出的是 $L_\infty$ 关系式（斜对称化之后的结果）；要得到 $A_\infty$ 关系式，必须把系数代数换成上三角矩阵代数 $\mathcal A$ 再张量，见注 3.2 与式 (47)。作者在 (B) 里没有点破这一层，读者若不留意会误以为 $A_\infty$ 关系式可以直接从 BV 主方程读出。\hfill{\footnotesize （源码第 54 行）}
  \item (A) 中“取奇异链”不是任意选择：$F_a-F_a\equiv0$ 的临界点集是整个 $X$，故自同态的 Morse 复数形式上就是 $X$ 上的链；取奇异链（而非只取点）使复合映射 (4) 能作用在“被看作分布”的链上。若改取 de Rham 形式或上链，会付出别的代价，见第 2.7 节与我在此处的批注。\hfill{\footnotesize （源码第 56 行）}
\end{enumerate}

\subsection*{第 2 节、第 2.1 节}
\begin{enumerate}
  \item 读本节前先记住一个总原则：这个范畴里所有“数出来的整数”都带符号，符号由 Morse 理论的有向计数给出；而所有横截性问题都靠“假设 2.2”回避，而不是靠扰动具体实现。作者把这一取舍称为“部分定义”，这是理解后文两个证明为何都用“无穷小邻域/光滑化”操作的前提。\hfill{\footnotesize （源码第 18 行）}
  \item 式 (1) 的下标约定极易出错：$P_{ab}$ 中前一个下标 $a$ 是源（$F_a$），后一个 $b$ 是靶（$F_b$），与函数差写成 $F_a-F_b$ 的顺序一致。若误按 $F_b-F_a$ 理解，Morse 指标会整体反号，式 (4) 的树形规则也随之颠倒。核对原书 p.4 式 (4) 与其后的三价顶点规则。\hfill{\footnotesize （源码第 40 行）}
  \item 提示“按余维数分级”与“按余指标分级”是同一件事的两种说法，但只差一个整体次数平移：$\operatorname{codim}=\dim X-\operatorname{coind}$。作者在 (1)(2) 中并列使用两种分级，正是为了让式 (3) 的次数公式 $2-n$ 在两块上都成立；这一处若统一成“按维数分级”，所有公式都要整体加 $\dim X$ 的平移，切勿随手统一。\hfill{\footnotesize （源码第 42 行）}
  \item 式 (4) 中 $\#\mathcal M$ 的参数写作 $P_{12}$，而式 (4) 左侧与求和处写作 $P_{a_1a_2}$。这是原文的记号不一致（疑为缩写或笔误），译文照录不改；读者把它读作 $P_{a_1a_2}$ 即可。\hfill{\footnotesize （源码第 88 行）}
  \item 三价顶点的颜色规则 $(a,b),(b,c)\to(a,c)$ 是说：沿树从叶到根走，标号必须“接力”，且每经过一个三价顶点就换一次色。这条规则把树的组合型与复合的输入顺序绑在一起，是全节最实质的约定；它同时也是“不需要四价顶点”的原因——四价顶点的两种分解在同调层面相消（见第 4.2 节 (a) 类边界层片的讨论）。\hfill{\footnotesize （源码第 90 行）}
  \item 这里“反推”用的是循环形式的配对 $\langle x_0,m(x_1,\dots,x_p)\rangle$，而不是 (4) 的非循环版本；两者相差一个取对偶的手续，符号来自 $\dim X$ 次的交换（见前一条批注）。若直接对 (4) 取 $a_1=a_r$，会多出一个不属于树的边界项，这是最容易算错的一步。\hfill{\footnotesize （源码第 106 行）}
  \item 式 (7) 中的 $\pm$ 号在原文里没有显式写出，它是 Koszul 符号：当内层 $m$ 的次数为奇时要对换序引入符号，具体为 $(-1)^{|x_1|+\cdots+|x_r|}$ 型的因子。这是全节唯一“不显式”的地方，定理 1 的两个证明（第 3、4 节）都给出同一符号约定；若读者自行补证，务必先固定 Koszul 约定，否则会差出整体符号。\hfill{\footnotesize （源码第 118 行）}
  \item 定理 1 的证明不是本节给出，而是第 3、4 节的两个独立证明。这是本文的组织方式：先把范畴定义“数”出来，再用两种场论把关系式“证”出来。阅读时建议先接受定理 1，直接进入第 3 节看有效作用量如何把它变成一个 BV 主方程。\hfill{\footnotesize （源码第 120 行）}
\end{enumerate}

\subsection*{第 2.2–2.3 节}
\begin{enumerate}
  \item 图 1 中边上的数字 1、2、3 是三个对象的标号（对应图 2 的 $L_1,L_2,L_3$ 即圆盘的三段边界弧），不是边的编号。每条边另带双指标 $(a,b)$ 与方向：$P_{12}$、$P_{23}$ 是叶，$P_{13}$ 是根，方向由叶指向根；方向读反，复合的输入与输出就对调。\hfill{\footnotesize （源码第 30 行）}
  \item 图 2 中间那幅画在 $T^*X$ 里：$L_a=\mathrm{graph}(\epsilon\,\mathrm{d}F_a)$ 是恰当 Lagrange 子流形，投影 $T^*X\to X$ 把它压到 $X$ 上。$\epsilon\to0$ 时全纯圆盘的面积趋于零，圆盘的极限是梯度流树；这里的“全纯”一律指伪全纯（见第 1 节脚注 1）。紧性论证原文未写，可参看 \cite{fukayaoh}。\hfill{\footnotesize （源码第 41 行）}
  \item 维数公式逐项对账：$\operatorname{ind}$ 是 $F_1-F_2$ 的 Morse 指标；$P\to P'$ 的轨道模空间维数为 $\operatorname{ind}(P)-\operatorname{ind}(P')$，再减 1 是去掉梯度流的 $\mathbb R$-平移。每个 2 价顶点是一条边与链 $Z$ 的交点：链的余维数为 $c$ 时只贡献 $c-1$ 个条件，因为交点沿 1 维的边可自动滑动。另注意 (1) 按余指标分级，此处用指标，二者只差常数 $\dim X$，差值不变。\hfill{\footnotesize （源码第 79 行）}
  \item 把 (7) 按 $n=2$ 展开：$r=0,s=1$ 给出 $m(\mathrm{d}x_1,x_2)$，$r=1,s=1$ 给出 $m(x_1,\mathrm{d}x_2)$，$r=0,s=2$ 给出 $\mathrm{d}m(x_1,x_2)$，正是 (8) 的三项。注意 $x_2=Z$ 是链，其微分是边界算子，故第二项写作 $m([P],\partial Z)$ 而非 $m([P],\mathrm{d}Z)$；$x_1=[P]$ 是 Morse 链，$\mathrm{d}_{\mathrm{Morse}}[P]=[R]$ 即式 (8) 中的下花括号。\hfill{\footnotesize （源码第 100 行）}
  \item (9) 中 $\mathcal L_v$ 是李导数，$\iota_v$ 是内乘（interior product）。$U$ 的 0 次部分是拉回 $e^{-t\mathcal L_v}$，1 次部分是 $-\mathrm{d}t\,e^{-t\mathcal L_v}\iota_v$，故 (12) 的负号来自这里；$\int_0^\infty U(t,\mathrm{d}t)$ 只取 $\mathrm{d}t$ 分量，即在 $(0,+\infty)$ 上积一个 1 次形式。把 $\iota_v$ 从 $e^{-t\mathcal L_v}$ 右边挪到左边不变号，因为 $v$ 在自身的流下不变：$e^{-t\mathcal L_v}\iota_v=\iota_v e^{-t\mathcal L_v}$。\hfill{\footnotesize （源码第 120 行）}
  \item $i$ 用不稳定流形、$p$ 用稳定流形，二者不可互换。$p\circ i=\mathrm{id}$ 的道理是：$\delta_{\mathrm{Unstab}_Q}$ 的余维数为 $\operatorname{ind}Q$，沿维数为 $\operatorname{ind}P$ 的 $\mathrm{Stab}_P$ 积分，只有 $\operatorname{ind}P=\operatorname{ind}Q$ 时可能非零；此时 $P\to Q$ 的梯度轨道必为常值，故只有 $P=Q$ 贡献 1。\hfill{\footnotesize （源码第 147 行）}
  \item $U=e^{-(\mathrm{d}_t+[\mathrm{d}_X,-])(t\iota_v)}$ 是简写，不是通常的算子指数：指数里的因子并不交换。用 $[\mathrm{d},\iota_v]=\mathcal L_v$ 可算出 $(\mathrm{d}_t+[\mathrm{d}_X,-])(t\iota_v)=t\mathcal L_v+\mathrm{d}t\,\iota_v$；作者实际取的定义是 (9)，两种写法给出同一个满足 (14) 的形式。\hfill{\footnotesize （源码第 164 行）}
  \item 原文 p.8 此处作 the gradient trajectory of starting at y，多一个 of，疑为笔误；译文按 the gradient trajectory starting at y 理解，其余照录。另注意 $\widetilde Y$ 不是子流形：它是映照 $(t,y)\mapsto(\mathrm{Flow}_t(v)(y),y)$ 在 $(0,+\infty)\times X$ 上的像，边界由 $t=0$ 处的 $\mathrm{Diag}$ 与 $t=+\infty$ 处的 $Y_{+\infty}$ 拼成，这里的 $\partial$ 是链（current）意义下的边界。\hfill{\footnotesize （源码第 172 行）}
  \item (18) 中间的等号关键是 $R_\epsilon$ 与恒等映射链同伦：$R_\epsilon-\mathrm{id}=\mathrm{d}h_\epsilon+h_\epsilon\mathrm{d}$，而 $\delta^\epsilon_{Z_1}$ 闭，被积式又是顶次形式。于是把 $\delta^\epsilon_{Z_2}$ 换回 $\delta_{Z_2}$ 只差一个恰当形式的积分，在闭流形 $X$ 上为零。原文只写“$\epsilon$ 充分小”，这一步省略了。\hfill{\footnotesize （源码第 192 行）}
  \item 正规化的动机：未正规化的 $i$ 取值在分布形式中（含 delta），而 $p$ 只定义在光滑形式上；后文的树形求和要反复做形式乘法，故必须先把对象换成光滑的。$K^\epsilon=R_\epsilon K+h_\epsilon$ 中的 $h_\epsilon$ 不是可选项：它恰好补上 $R_\epsilon$ 与 $\mathrm{id}$ 之间的链同伦缺口，缺了它 (21) 不成立。\hfill{\footnotesize （源码第 224 行）}
\end{enumerate}

\subsection*{第 2.4–2.6 节}
\begin{enumerate}
  \item 此处极易把符号搞反：正确的约定是 $|T_\alpha|=1-\operatorname{codim}C_\alpha$，而不是 $\operatorname{codim}C_\alpha-1$。按前者，例 4 中 $\operatorname{codim}C=1$ 给出 $|T|=0$，故 $e^T$ 的展开不带附加符号；若按后者则 $|T|=-2$，$m(T)$ 的奇偶性与 $e^T-1$ 的符号会整体改变，式 (25) 的符号也随之失配。\hfill{\footnotesize （源码第 27 行）}
  \item 为什么模型中要张量 $\mathbb Q$：Morse 链复形本来定义在 $\mathbb Z$ 上，但重复闭链的位移正规化（注~\ref{rem:2.1}）会产生 $1/k!$ 型的有理系数，$m(T)$ 因而只在张量 $\mathbb Q$ 之后才成为良定义的自同态。若忽略这一点，会误以为 $m(T)$ 保持整系数，进而误判整系数版本是否存在。\hfill{\footnotesize （源码第 35 行）}
  \item 对称化给出的权重是 $T^2/2$，不是笔误：$C$ 的两份副本各得 $1/2$ 的权重，分母上的 $2=2!$；推广到 $k$ 份相互重复的副本时权重为 $1/k!$，这正是例 4 与本节第二个脚注（原书脚注 8）中系数 $1/k!$ 的来源。只有当位移足够小时，对称化的结果才等于 $\tfrac12\sum_\gamma(\#(\gamma\cap C))^2$。\hfill{\footnotesize （源码第 43 行）}
  \item 为什么 $\mathrm{d}_{\mathrm{Morse}}=0$：在圆上，从极大点 $P$ 到极小点 $Q$ 恰有两条梯度轨道（沿两侧各一段半圆），二者定向相反，带符号计数相互抵消。而 $m_P^Q(T)$ 数的是经过闭链副本的轨道：对每组位移，按正确次序经过副本的轨道恰有一条，对称化后得 $1/k!$，累加即 $e^T-1=\sum_{k\ge 1}T^k/k!$。\hfill{\footnotesize （源码第 71 行）}
  \item 作者说“显然”，实际用到两步：其一，$\epsilon\to0$ 时 $\delta^\epsilon_C$ 弱收敛到 $\delta_C$，$K^\epsilon$ 的积分核收敛到 $\delta_{\widetilde Y}$（式 (17) 与第 2.3 节的正则化构造）；其二，沿 $X^V$ 的积分在横截性下化为交点计数。标准出处是 Gugenheim \cite{gugenheim} 与 Lambe--Stasheff \cite{lambe} 的扰动引理，综述见 Crainic \cite{crainic}。\hfill{\footnotesize （源码第 104 行）}
  \item 本段 $m=r-2+\sum_i k_i$ 中的 $m$ 是空间 $X^m$ 的指数（内部顶点总数），与复合映射 $m$ 无关；同段 $Z_{\alpha,a}$ 的下标 $a$ 与临界点 $P_{ab}$ 的双指标也各有含义，读式 (28) 时先分清这两套指标。\hfill{\footnotesize （源码第 132 行）}
  \item $\pi_e^{-1}\widetilde Y_e$ 里的波浪号不能丢：$\widetilde Y_e=\bigcup_{t>0}Y_t$ 是式 (17) 中 $X\times X$ 内的整条扩展轨道族的链，而 $Y_t$ 只是 $t$ 时刻的流动像。$\pi_e$ 把 $X^V$ 投到第 $q_1,q_2$ 份 $X\times X$，$\pi_q$ 只投到单个 $X$，两者的纤维维数不同，把 $\pi_q^{-1}(\widetilde Y_e)$ 写进公式是错的。\hfill{\footnotesize （源码第 143 行）}
  \item 疑似笔误：原书 p.13 把维数公式排成两行，第一行以 $\dim\mathcal M=$ 结尾，第二行又以 $=$ 起头，于是等号出现两次。译文照印，未删去第二个等号；按上下文，第一个等号后应直接接第二行的右端。\hfill{\footnotesize （源码第 161 行）}
  \item “预期维数为零”来自式 (28) 后的维数公式，作者没有写出代入过程：此处 $\operatorname{coind}(\mathbf S)-1=2-1=1$，$-\bigl(\operatorname{coind}(\mathbf N)-1\bigr)=-(0-1)=1$，两条链贡献 $-\bigl((1-1)+(2-1)\bigr)=-1$，最后再减 1，合计 $1+1-1-1=0$。实际模空间是一维的，故横截性必然失效。\hfill{\footnotesize （源码第 182 行）}
  \item 疑似笔误：式 (31) 的乘积写作 $\prod_{a=1}^{N-1}$，但被积函数里是 $\delta_{\mathrm{Unstab}_{P_{a_ia_{i+1}}}}$，下标是 $i$ 而不是 $a$；照式 (29) 应为 $\prod_{i=1}^{r-1}$，或把被积函数的下标改成 $P_{a\,a+1}$。另外这里的 $N-1$ 与式 (29) 的 $r-1$ 也不一致（$N$ 是对象个数，$r$ 是树中叶子个数）。译文照原样排印，未擅改。\hfill{\footnotesize （源码第 233 行）}
  \item 末句的“波前集横截”是关键：Hörmander 的波前集理论断言，两个分布之积有定义当且仅当二者的波前集沿对角逐点横截。$\delta_Z$ 的波前集是 $Z\subset X$ 的余法丛，故对 $X^V$ 上各拉回形式的乘积而言，横截性条件恰是假设~\ref{ass:2.2} 的相交横截性。可参看 Hörmander《The Analysis of Linear Partial Differential Operators》第 8 章。\hfill{\footnotesize （源码第 235 行）}
\end{enumerate}

\subsection*{第 2.7 节、第 3 节}
\begin{enumerate}
  \item 这里三类 $\delta$ 的含义不同：$\delta_{Z_q}$ 是链的 delta 形式（对链积分即数交点），$\delta_{\mathrm{Unstab}_{P_{a_ia_{i+1}}}}$ 与 $\delta_{\mathrm{Stab}_{P_{a_1a_r}}}$ 是稳定、不稳定流形的 delta 形式，而 $\epsilon$-光滑化 delta 形式是链的 delta 形式的正则化。把三类混为一谈，就会误以为链与临界点在这个变体里地位相同。\hfill{\footnotesize （源码第 16 行）}
  \item 为什么 cup 积不需要横截性、而链的相交需要？cup 积是上链层面的代数运算，任意两条上链都可以相乘；链的相交则要求子流形横截，否则交点集不是离散的，计数无意义。但高阶复合仍要数 Morse 树与链的交，所以横截性没有消失，只是从基本运算挪到了复合这一步——这正是作者说「仍得到部分定义的范畴」的原因。\hfill{\footnotesize （源码第 20 行）}
  \item 本节与第 4 节是定理~\ref{thm:1} 的两个独立证明。核心工具是同伦转移的存在性：给定收缩数据 $(i,\mathrm{p},K)$ 后，只要树形求和按次数逐次有限（Morse 收缩经正规化后满足这一条），dg 代数结构就能搬到 $\mathbb{M}$ 上。标准陈述是同调扰动引理，见 \cite{crainic,gugenheim,lambe}；树形求和的具体形状即 Kontsevich--Soibelman 公式 \cite{kontsevichsoibelman}。\hfill{\footnotesize （源码第 32 行）}
  \item 收缩三元组的标准条件是 $\mathrm{p}\,\mathrm{i}=\mathrm{id}_{\mathbb{M}}$ 与 $\mathrm{i}\,\mathrm{p}-\mathrm{id}_V=d_VK+Kd_V$（符号随分级约定略有出入），且 $K$ 的次数为 $-1$。$(a,a)$ 分块取 $\mathrm{i}=\mathrm{p}=\mathrm{id}$、$K=0$，是全篇「增强」的来源：$F_a-F_a\equiv0$ 没有 Morse 理论可用，故 $\Omega_{aa}$ 原样留在 $\mathbb{M}$ 中，不做 Morse 收缩。\hfill{\footnotesize （源码第 66 行）}
  \item 次数核对：$n$ 叶二叉有根树有 $n-1$ 个内部顶点、$n-2$ 条内部边；乘积次数为 $0$，$K$ 的次数为 $-1$，故 $\mu_n$ 的次数为 $2-n$，与式~\eqref{eq:3} 中复合映射的次数一致。若把 $K$ 与乘积的次数记反，式~\eqref{eq:36} 的符号会整体错位。\hfill{\footnotesize （源码第 83 行）}
  \item 上式原书是手绘的二叉树（叶装饰 $\mathrm{i}(X_1),\mathrm{i}(X_2),\mathrm{i}(X_3)$，先由乘积合并前两个叶，在内部边上作用 $K$ 后再与第三个叶相乘，最后在根处作用 $\mathrm{p}$）。译文按原式整行裁图插入（图文件 \mbox{figtree33.png}），未重排；其组合型由「$n$ 叶二叉有根树」与上面的装饰规则唯一决定。\hfill{\footnotesize （源码第 84 行）}
  \item $X_i$ 落在哪个分块极易记错：$a_i\neq a_{i+1}$ 时 $X_i=x_i$ 是 $\mathbb{M}_{a_ia_{i+1}}=\mathrm{MC}(F_{a_i}-F_{a_{i+1}})$ 中的临界点或 Morse 链；$a_i=a_{i+1}$ 时 $X_i=\delta^{\epsilon}_{x_i}$ 落在 $\Omega_{a_ia_i}$ 中。$\delta^{\epsilon}_{x_i}$ 不是临界点，必须当作光滑化 delta 形式才能进入树形公式~\eqref{eq:33}。\hfill{\footnotesize （源码第 96 行）}
  \item 这一步是式~\eqref{eq:33} 与式~\eqref{eq:31} 逐项对应的组合核心。式~\eqref{eq:31} 的树允许 2 价顶点来承载链 $Z_{\alpha,a}$；而式~\eqref{eq:33} 的二叉树叶在 $a_i=a_{i+1}$ 时带的是 $\delta^{\epsilon}_{x_i}$。截去这样的叶及其边，相当于把它并进 $\delta_{Z}$ 型因子，同时在 $T'$ 中留出一个 2 价顶点。图~\ref{fig:8} 给出的正是最简单的例子。\hfill{\footnotesize （源码第 102 行）}
  \item 原书 p.17 把 Soibelman 拼作 Soibleman，疑为笔误；译文照录，未擅改。读者应把它与 \cite{kontsevichsoibelman} 中的 Soibelman 认作同一人，切勿当成另一文献。\hfill{\footnotesize （源码第 118 行）}
  \item 式~\eqref{eq:37} 有三层结构：$A$ 的矩阵指标在 $N\times N$ 上（对应 $\mathbb{F}$ 的 $N$ 个对象），每个矩阵元再取值于 $\mathcal A$（大型上三角矩阵代数）。$\langle\,,\rangle$ 是 $\mathfrak g$ 与 $\mathfrak g^*$ 的配对；对微分形式还要配上 wedge 积与矩阵迹，原文按标准 BF 记号省略未写。$A$ 是奇变量，故 $dA+\frac12[A,A]$ 就是曲率。\hfill{\footnotesize （源码第 136 行）}
  \item 为何必须把系数扩成 $\mathcal A$：只取 $N\times N$ 矩阵时，反对称化后只能读到 $L_\infty$ 括号，$\Ainf$ 的 $\mu_n$ 被混同。张量 $\mathcal A$ 之后，矩阵单位 $t_{ij}$ 像「标签」一样把输入与输出的次序记下来（见式~\eqref{eq:47}）。再取正向极限 $\widetilde N\to\infty$，是为了不设 arity 的上限。$\mathcal A$ 幂零这一点，在「非幂零 $\mathfrak g$」那条脚注处会用到。\hfill{\footnotesize （源码第 142 行）}
  \item 式~\eqref{eq:38} 把「几何」塞进了规范固定：$\iota_{v_{ab}}$ 是沿梯度场 $v_{ab}$ 的缩并。作者断言这个条件等价于「$(A_W)_{ab}$ 被 Morse 链同伦式~\eqref{eq:12} 湮灭」。结合下文 $\mathcal L=\mathrm{im}K\oplus\mathrm{im}K$，这条等价关系正是「几何规范固定转化为内部边装饰 $K$」的翻译词典，也是全节物理思想的关键一步。\hfill{\footnotesize （源码第 150 行）}
  \item 式~\eqref{eq:40} 两端的方括号是次数平移：$[1]$ 与 $[\dim X-2]$ 必须逐项对齐，混淆两者会破坏奇辛配对的次数。右端 $\mathbb{M}\otimes\mathcal A^*$ 中的星号是对 $\mathcal A$ 取对偶，不是对整个张量积取对偶。$\mathcal L=\mathrm{im}K\oplus\mathrm{im}K$ 说明纤维积分沿 $K$ 的像进行，故内部边最终由 $K$ 装饰。\hfill{\footnotesize （源码第 166 行）}
  \item 系数与 $|\mathrm{Aut}(\Gamma)|$ 对应：前三个单链树的权重是 $1/2$；最后的「哑铃」树两条 $K$ 边可交换，自同构群阶为 $8$，故系数是 $1/8$ 而不是 $1/4$。抄写式~\eqref{eq:42} 时最易在此出错，请对原书 p.18 第二行逐字核对。\hfill{\footnotesize （源码第 191 行）}
  \item 六项分工要记清：(b)(c) 只含 $A_W$，给出传播子与三价相互作用；(a)(e) 不含 $A_W$，在纤维积分中是常数；(d)(f) 各带一条 $A_Z$ 或 $B_Z$ 的外线，把慢场接到动力学部分。注意 (e) 与 (f) 的被积式都含 $\frac12[\cdot,\cdot]$，但 (e) 是常数项、(f) 是外线顶角，切勿记反。\hfill{\footnotesize （源码第 221 行）}
  \item $1/n!$ 来自 $n$ 个等价输入 $A'$ 的置换（$l_n$ 斜对称，$n!$ 个置换给出同一项）。按叶数归并式~\eqref{eq:41} 时，这个因子必须与树的同构群阶数一起核对，否则会差整数倍。\hfill{\footnotesize （源码第 243 行）}
  \item 这一段是全节的物理枢纽，值得慢读：规范固定（几何条件）→ 传播子（Feynman 规则里的边装饰）→ Morse 链同伦 $K$（组合数据）→ Morse 树计数（范畴的复合）。四个对象其实是同一件事的四种语言，作者用「规范理论的逻辑」把它们串成一条链。若省略这一步，读者会误以为 BF 有效作用量与 Morse 树计数只是碰巧形式相似。\hfill{\footnotesize （源码第 247 行）}
  \item 式~\eqref{eq:46} 与式~\eqref{eq:36} 的求和范围不同：这里 $r+s=n$，$n$ 个输入由置换 $\sigma$ 分配进内外两个括号，全部参与；而 $\Ainf$ 的式~\eqref{eq:36} 允许 $r+s\le n$，尾部可有剩余输入。照抄式~\eqref{eq:36} 的范围会漏项。\hfill{\footnotesize （源码第 260 行）}
  \item $t_{ij}$ 是上三角矩阵代数 $\mathcal A$ 中的矩阵单位。取 $t_{12},t_{23},\dots,t_{n,n+1}$ 是沿态射链 $(a_1,a_2),\dots,(a_n,a_{n+1})$ 把指标串起来，输出恰好落在 $t_{1,n+1}$ 上。若 $b_i\neq a_{i+1}$，乘积在 $\mathcal A$ 中为零，式~\eqref{eq:47} 两端同时消失。这正是 $\mathcal A$ 能「记住」$\Ainf$ 复合的组合机制。\hfill{\footnotesize （源码第 268 行）}
\end{enumerate}

\subsection*{第 4 节}
\begin{enumerate}
  \item 本节是第~\ref{sec:3} 节的“对偶”版本。第~\ref{sec:3} 节走二次量子化：把 $X$ 上的 BF 理论作纤维 BV 积分，树是 Feynman 图。本节走一次量子化：树本身是单粒子超量子力学的“时间轴分叉”，模空间 $MT$ 只负责记录各边长度。两条路线给出同一批复合映射，但本节只需 $MT$ 上一个闭形式加 Stokes 定理。\hfill{\footnotesize （源码第 15 行）}
  \item 式 (48) 与式 (9) 是同一个算子：式 (9) 把它写成 $e^{-tL_v}(1-\mathrm{d}t\,\iota_v)$。因为 $(\mathrm{d}t\,\iota_v)^2=0$ 且 $[L_v,\iota_v]=0$，指数可以拆开，两种写法一致。关键在于每次出现 $\mathrm{d}t$ 必同时出现一个 $\iota_v$，这条“捆绑”是全节次数计数的起点。\hfill{\footnotesize （源码第 46 行）}
  \item 这条“闭形式加边界分解”的论证与上同调场论中 WDVV 方程的证明同构：那里模空间是带四个标记点的球面模空间（复一维），边界是三个点；这里模空间是一维的树模空间，边界只由 $t\to0$ 与 $t\to+\infty$ 组成。两者都归结为一组二次关系，一维版本就是 $\Ainf$ 关系式。见 \cite{losev,kontsevichmanin}。\hfill{\footnotesize （源码第 52 行）}
  \item “内部边”不含紧挨 1 价顶点的边，那类边长度恒为 $+\infty$。因此 $\dim MT$ 等于内部边条数，图~\ref{fig:11} 的三棵树各有 $t$、$t'$、$t''$ 三条内部边，都是 3 维胞腔。$t\to0$ 是 UV 边界（收缩），$t\to+\infty$ 是 IR 边界（分解），两者在式~\eqref{eq:55} 的左端分别对应 (a)–(d) 类与 (e) 类层片。\hfill{\footnotesize （源码第 72 行）}
  \item 原文 (49) 在两行的行末与行首各印了一个 $\otimes$，照排后出现两个张量积符号，这是排版残留，译文照印未改。另注意输入顺序与式 (3) 一致：先 $\Mor_{a_1,a_1}^{\otimes k_1}$、再 $\Mor_{a_1,a_2}$，例 6 的排法即由此而定；次序写反则 $\int_{MT}I=m$ 不成立。\hfill{\footnotesize （源码第 90 行）}
  \item (50) 的次数核对法：$U=e^{-tH}(1-\mathrm{d}t\,G)$ 中 $\mathrm{d}t$ 与 $G=\iota_v$ 成对出现，每对净次数为 0，但使 $X$ 方向的形式次数降 1。要对 $X$ 与 $MT$ 同时取顶次项，需三条内部边各取一个 $\mathrm{d}t$（凑出 $\mathrm{d}t\wedge\mathrm{d}t'\wedge\mathrm{d}t''$），且 $\delta_{\mathrm{Stab}P}$、$\delta_{\mathrm{Unstab}P}$、$\delta_Z$ 的次数之和（分别为 $P$ 的指标、$P$ 的余指标、链的余维数）等于 $\dim X+3$。\hfill{\footnotesize （源码第 110 行）}
  \item 原文此处印作 “the respective states of states”，疑为 spaces of states 之笔误。译文照印并加引号，读者按“态空间”理解即可。之所以要光滑化：$\delta_{\mathrm{Unstab}P}$ 是分布形式，两个分布形式相乘没有定义，换成 $\delta^\epsilon$ 后所有运算都在光滑形式内完成，最后再取 $\epsilon\to0$。\hfill{\footnotesize （源码第 116 行）}
  \item 这里的 $Q$ 要读成作用在算符上的分次交换子 $[\mathrm{d}_X,\cdot]$：在微分形式上它就是外微分 $\mathrm{d}$（即第 4.1 节开头的 $Q=\mathrm{d}$），在算符上则是 $[Q,\cdot]$。式 (14) 的 $(\mathrm{d}_t+[\mathrm{d}_X,-])U=0$ 与 (51) 是同一件事。\hfill{\footnotesize （源码第 127 行）}
  \item 式 (52) 的求和条件用位置指标 $i,j$（$k'_i+k''_i=k_i$），式 (53) 却改用 $a,b$（$k'_a+k''_a=k_a$）；按上下文二者指同一对指标，原文如此，译文照印。另注意 (52) 第一因子的颜色列表 $a_1,\dots,a_j$ 有 $j$ 项，而 $k$-序列 $k'_i,k_{i+1},\dots,k_{j-1},k'_j$ 只有 $j-i+1$ 项，项数不符；按剪开 $(a_i,a_j)$-边的几何，第一因子应接 $a_i,\dots,a_j$。原文如此，译文照印，未擅改。\hfill{\footnotesize （源码第 145 行）}
  \item 式 (53) 的“立即推出”指：把式 (15) 的 $U(+\infty)=i\circ p$ 作用到那条趋于无穷的边上，$I$ 就裂成两个因子，而 $i\circ p$ 正好把一侧的输出与另一侧的输入对接。对每条趋于无穷的边归纳一次即得 (53)，作者省去了这个逐边归纳。\hfill{\footnotesize （源码第 146 行）}
  \item 式 (55) 的 $\pm$ 只来自边界定向的约定：由 (51) 得 $\mathrm{d}_{MT}I=-QI$，两边在 $MT$ 上积分并用 Stokes 定理，左端化为 $\int_{\partial MT}I$，右端为 $-Q\int_{MT}I$（$Q$ 作用在 $X$ 一侧，与 $MT$ 上的积分可交换）。边界取外法向或内法向定向，差别恰是整体符号，故原文写 $\pm$。\hfill{\footnotesize （源码第 161 行）}
  \item 式 (56) 里 $m$ 有三种用法，极易混：$r=2$ 的 (4)（如 $m(Z_1,Z'_1,[P])$）、$r=1$ 的 (5)（如 $Z_1\cap Z'_1$ 与 $\partial Z_1$）、以及外层 $m$ 吞掉一个内层 $m$ 的嵌套（如 $m(m(Z_1,Z'_1,[P]),Z_2)$）。核对时先按输入的颜色序列判断该用 (4) 还是 (5)。\hfill{\footnotesize （源码第 188 行）}
\end{enumerate}

\subsection*{第 5 节}
\begin{enumerate}
  \item 余维一闭链与平坦联络的对应来自式 (5)：$m(Z_1,Z_2)=Z_1\cap Z_2$ 是链的带符号相交。取 $Z$ 为余维一闭链时，相交数给出绕该闭链一周的相位，这正是和乐；把系数换成 $GL(\nu_a)$ 矩阵，和乐就成为非阿贝尔的。作者没有展开这个对应，正文只用了「可以解释为」的说法。\hfill{\footnotesize （源码第 16 行）}
  \item “堆叠”是 $T^*X$ 的 Fukaya 范畴中的标准图像：一个 Lagrange 子流形取 $\nu$ 个重合副本，其上的规范场联络自然取 $GL(\nu)$ 值。注意这里的副本不是“扰动开来的”，而是重合的；若真的扰动开，全纯圆盘会重新出现，反而不属于本文的框架。\hfill{\footnotesize （源码第 18 行）}
  \item 末句把 §2.1 的“特设”模型解释成一般构造的推论：$F_a-F_a\equiv0$ 的临界集是整个 $X$，故自同态的 Morse--Bott 链复形就是 $X$ 上的奇异链，与式 (2) 一致。但要注意 Morse--Bott 微分比链的边界算子多一项——临界流形之间的梯度流；在 $a=b$ 的特例里这一项退化，所以 (2) 只需边界算子。\hfill{\footnotesize （源码第 24 行）}
\end{enumerate}

\endinput
